% Décharge d'un condensateur. Résolution analytique
\begin{minipage}{0.55\linewidth}
    Un condensateur initialement chargé est inséré dans le montage suivant~:

    À la date $t=0$, la charge du condensateur a pour valeur~:
    $q_{A}=\qty{0.6e-6}{\C}$ et on ferme l'interrupteur $\mathrm{K}$.
    \begin{enumerate}[series=y]
        \item Représenter par une flèche le sens de circulation du courant
              d'intensité $i$ dans le circuit.
        
              Représenter par des flèches les tensions $u_{C}$ aux bornes du
              condensateur et $u_{R}$ aux bornes du conducteur ohmique, en
              respectant la convention récepteur.
        \item Établir une relation~(1) entre $u_{R}$ et $u_{C}$.
    \end{enumerate}
\end{minipage}
\hfill
\begin{minipage}{0.44\linewidth}
    \begin{figure}[H]
        \centering
        \begin{circuitikz}
            \newdimen\d \d=3cm
            \draw (0,0) coordinate (O)
                to[short] ++(0,\d)
                to[cute open switch,l=$\mathrm{K}$] ++(\d,0)
                to[R,l=$R$] ++(\d,0)
                to[short,-*] ++(0,-\d) node[below] {$B$}
                to[C,l_=$C$,-*] (O) node[below] {$A$};
        \end{circuitikz}
    \end{figure}
\end{minipage}

\begin{enumerate}[resume=y]
    \item À l'aide de la relation~(1), établir l'équation différentielle à
          laquelle obéit $u_{C}$.
    \item Une solution de cette équation différentielle est de la forme~:
          $u_{C}=a\times\mathrm{e}^{bt}$. Déterminer les constantes $a$ et
          $b$. En déduire l'expression de $u_{C}$ en fonction du temps.
\end{enumerate}

\begin{donnees}
    \begin{itemize}
      \item $R=\qty{15}{\kilo\ohm}$ et $C=\qty{0.1}{\uF}$.
    \end{itemize}
\end{donnees}


